\scanpage{53}
\mspage{674}

the image of $\underline J$ in it is an ideal ...

We do find a ring structure on $\underline A_{[N]}$, and an exact sequence \deleted{\unsure{of}}
\[
\tag{71}
\diagram{\begin{tikzcd}[column sep=large]
1 \arrow[r,no head] & \underline J \arrow[r,no head] & \underline A_{[N]} \arrow[r] & B_{S_0} \arrow[r] & 1
\end{tikzcd}}
\]
where $\underline A_{[N]}\mathbin{-}B_{S_0}$ is a \unsure{morphism} of \uline{ring schemes}. Passing to the multiplicative groups \deleted{\unsure{rings}}, we find an exact sequence
\[
\tag{72}
\diagram{\begin{tikzcd}[column sep=large,row sep=small]
1 \arrow[r] & (1+\underline J)^* \arrow[r,no head] & A_{[N]}^* \arrow[r,no head] & (B_{S_0})^* \arrow[r,no head] \arrow[d,equal] & 1 \\
& & & (B^*)_{S_0} &
\end{tikzcd}}
\]
where for $A_{[N]}^*$, $(B_{S_0})^{\unsure{*|'}}$ and $B^*$, the $*$ denotes the \unsure{passage} to the multiplicative \unsure{subgroup} and where $(1+\underline J)^*$ has the analogous meaning
\[
\tag{73}
\diagram{\begin{tikzcd}[row sep=large]
(1+\underline J)^*\unsure{\subset|\subseteq}1+\underline J \arrow[d,draw=none,"="{description,rotate=90}] \\
\{(1+\lambda)\mid 1+\lambda\ \text{invertible in }\underline A\} \arrow[d,draw=none,"\simeq"{description,rotate=90}] \\
\{(\lambda,\lambda')\in\underline J\times\underline J\mid\lambda+\lambda'+\lambda\lambda'=0\ \text{in }\underline J\}.
\end{tikzcd}}
\]

Thus we will have, in the special case $A=\mathbb Z$, $J=N\mathbb Z$
\[
\tag{74}
\diagram{\begin{tikzcd}[column sep=large]
1 \arrow[r] & (\underline{1+N\mathbb Z})^* \arrow[r] & \mathbb O_N^* \arrow[r,no head] & (\mathbb Z/N\mathbb Z)^*_{S_0} \arrow[r,no head] & 1
\end{tikzcd}}
\]
where the kernel

\margin{52}
\margin{52\\\texttt{https://grothendieck.umontpellier.fr}\\Université de Montpellier, UMR CNRS 5149, Case courrier 051, Place Eugène Bataillon\\34095 Montpellier cedex 5 - France}

\scanpage{54}
\deleted{\mspage{\unsure{674}}}

$(1+N\mathbb Z)^*$ is defined as
\[
\tag{75}
(1+N\mathbb Z)^*\overset{\mathrm{dfn}}{=}
\{\lambda\in\deleted{\unsure{\mathbf G_a}}\underline{\mathbb Z}=\mathbf G_a\mid 1+N\lambda\ \text{invertible}\}
\]

\deleted{We have a canonical \unsure{homomorphism} (\unsure{in part}}

The $\mathbb O_N$ for varying $N$ form a projective system of Rings, the $\mathbb O_N^*$ form a projective system of commutative group schemes.

I am not going to look more closely at the structure of these projective systems, nor examine \unsure{the} projective limits --- \unsure{questions} which arise more \unsure{general|generally} when one starts from an $A$ as above. We are going to make explicit the canonical \unsure{homomorphism}
\[
\tag{76}
\diagram{\begin{tikzcd}[column sep=large]
\mathbb O_N^* \arrow[r,no head] & \mathbb O_1^*=\mathbb O^*=\mathbf G_m
\end{tikzcd}}
\]
\deleted{$\unsure{\varphi}$}

\deleted{whose kernel is}
which in general is neither an epimorphism, nor a monomorphism. Its kernel is \deleted{formed of the} \deleted{$\unsure{\varphi}\ \unsure{\lambda}$} isomorphic to the \unsure{subgroup|system} of
\unsure{$+$}
\[
\tag{77}
\{1+N[\lambda]\mid 1+N\lambda=1\ \text{i.e.}\ N\lambda=0\}
\]
which as a scheme (\unsure{without} group structure) is isomorphic to ${}_N\mathbf G_a$,

\margin{53}
\margin{53\\\texttt{https://grothendieck.umontpellier.fr}\\Université de Montpellier, UMR CNRS 5149, Case courrier 051, Place Eugène Bataillon\\34095 Montpellier cedex 5 - France}

\scanpage{55}
\mspage{675}

which is a group scheme (not flat if $N\ne1$) which is zero over $\operatorname{Spec}\mathbb Z[1/N]$, and \deleted{is \unsure{dimension|of dimension} relative $1$} the fibers at the $p\mid N$ are $\simeq\mathbf G_a$.

I realize now that it is a clumsy move to leave \unsure{the} \unsure{group} $\mathbb O_N^*$, which is \unsure{perfect} of its kind, to refer to $\mathbf G_m$.

I am going to examine in particular
\[
\tag{78}
\diagram{\begin{tikzcd}[column sep=large]
1 \arrow[r,no head] & (1+12\mathbb Z)^* \arrow[r,no head] & \mathbb O_{12}^* \arrow[r] & (\mathbb Z/12\mathbb Z)^*_{S_0} \arrow[r,no head] & 1.
\end{tikzcd}}
\]

Let us note that \deleted{$\mathbb Z/12\mathbb Z$}
\[
(\mathbb Z/12\mathbb Z)^*\simeq(\mathbb Z/4\mathbb Z)^*\times(\mathbb Z/6\mathbb Z)^*\simeq\mu\times\mu.
\]

Thus (78) takes the form
\deleted{and we find a set-theoretic section \unsure{the|the}\ \illegible{1}\ \unsure{of|of the}\ \unsure{group|kernel}\ \unsure{over}\ of the extension $\unsure{\beta}$}
\[
\tag{79}
\diagram{\begin{tikzcd}[column sep=large]
1 \arrow[r,no head] & (1+12\mathbb Z)^* \arrow[r,no head] & \mathbb O_{12}^* \arrow[r,no head,"\varepsilon_2\times\varepsilon_3"] & \mu\times\mu \arrow[r,no head] & 1
\end{tikzcd}}
\]

The quantities $\nu_{\unsure{2|3}}\ \nu_3$ are also \inserted{morphisms} defined on $\mathbb O_{12}^*$
\margin{\deleted{\unsure{here}}}
\[
\tag{80}
\diagram{\begin{tikzcd}[column sep=large]
\mathbb O_{12}^* \arrow[r,"(\nu_2,\nu_3)"] & \{0,1\}_{S_0}\times\{0,1\}_{S_0}
\end{tikzcd}}
\]
they are the composites
\[
\tag{81}
\diagram{$\begin{tikzcd}[ampersand replacement=\&,column sep=large,row sep=small]
\mathbb O_{12}^* \arrow[r,"(\varepsilon_2,\varepsilon_3)"] \& \mu\times\mu \arrow[r,"\sim"] \& \{0,1\}_{S_0}\times\{0,1\}_{S_0} \\
\& \& \text{deduced from}\quad\deleted{\unsure{(+1)\mapsto(0)}}\quad\begin{array}{rcl}1&\mapsto&0\\-1&\mapsto&1\end{array}
\end{tikzcd}$}
\]

\margin{54}
\margin{54\\\texttt{https://grothendieck.umontpellier.fr}\\Université de Montpellier, UMR CNRS 5149, Case courrier 051, Place Eugène Bataillon\\34095 Montpellier cedex 5 - France}

\scanpage{56}

\deleted{\unsure{that}}

and they are group homomorphisms (for addition modulo $2$ on $\{0,1\}_{S_0}$.

\inserted{We also have $\varepsilon_2=(-\unsure{1})^{\nu_2}=1-2\nu_2,\quad\varepsilon_3=(-\unsure{1})^{\nu_3}=1-2\nu_3$}

The ``four cases of \uline{signs}'' in $(\mathbb Z/12\mathbb Z)^*$ are given in the table
\[
\tag{82}
\diagram{$\begin{array}{r|c|c}
{}_{\varepsilon_3}\!\diagdown^{\varepsilon_2}
&\begin{array}{c}+1\\\text{i.e. }\nu_2=0\end{array}
&\begin{array}{c}-1\\\text{i.e. }\nu_2=1\end{array}\\ \hline
\begin{array}{r}+1\\\text{i.e. }\nu_3=0\end{array}&1&+5\\ \hline
\begin{array}{r}-1\\\text{i.e. }\nu_3=1\end{array}&-5&-1\\ \hline
\end{array}$}
\]

thus we find a set-theoretic section of (79) by taking the sections of $\underline{\mathbb O}_{12}^*$ of the form
\[
\tag{83}
\underline{\mu}_0=(-5)^{\nu_2}(+5)^{\nu_3}
\left\{\begin{array}{ll}
\equiv(-1)^{\nu_3}&(6)\\
\equiv(-1)^{\nu_2}&(4)
\end{array}\right.
\]

Thus a general section of $\mathbb O_{12}^*$ is written
\[
\tag{84}
\underline{\mu}=\underbrace{(-5)^{\nu_2}(+5)^{\nu_3}}_{\mu_0}+\unsure{12}[\unsure{\deleted{\lambda_0}\inserted{\nu_0}|\nu_0}]
\]
where
\[
\left\{\begin{array}{l}
(\nu_2,\nu_3)\in\Gamma(\{0,1\}_{S_0}\times\{0,1\}_{S_0})\\
\lambda_0\in\Gamma(\mathbf G_a).
\end{array}\right.
\]

\deleted{We \unsure{argue}\ \unsure{by|more}\ \unsure{modulo}\ $4$}
\inserted{Let us note that we have}
\[
\deleted{\mu=(4+1}
\qquad
5^{\nu_3}=1+4\nu_3
\]
and
\[
(-5)^{\nu_2}=\left[\begin{aligned}
(-1)^{\nu_2}5^{\nu_2}&=(1-2\nu_2)(1+4\nu_2)\\
&=1+2\nu_2-8\nu_2^2
\end{aligned}\right]=1-6\nu_2
\]
\margin{\uline{N.B.} $\nu_2=\nu_2^2$}

thus
\[
\tag{\unsure{85}}
\begin{aligned}
\underline{\mu}_0&=(1+4\nu_3)(1\unsure{-6}\nu_2)\deleted{\unsure{-8}\nu_2^2)}\\
&\deleted{=1+2\nu_2+\unsure{(-8)}\nu_2^2+4\nu_3(1+2\nu_2-8\nu_2^2)}\\
&\deleted{=1\unsure{-}2\nu_2+4[\unsure{\nu_2}-2\nu_2^2+\nu_3(1}\\
&=1+4\nu_3-6\nu_2-24\nu_2\nu_3
\end{aligned}
\]

\deleted{We \unsure{could} take}
\[
\deleted{(\unsure{85'|85})\qquad\mu'_0=1+4\nu_3-6\nu_2}
\]

\margin{55}
\margin{55\\\texttt{https://grothendieck.umontpellier.fr}\\Université de Montpellier, UMR CNRS 5149, Case courrier 051, Place Eugène Bataillon\\34095 Montpellier cedex 5 - France}

\scanpage{57}
\mspage{676}

If we \unsure{have} ((84) and (85))
\[
\underline{\mu}=\underline{\mu}_0+12[\unsure{\deleted{\lambda_0}\inserted{\nu_0}|\nu_0}]
=1+4\nu_3-6\nu_2+12(\nu_0-2\nu_2\nu_3)
\]
we conclude from this (by \unsure{considering} the congruences modulo $4$)
\[
\tag{86}
\mu=1-2\nu_2+4\underbrace{[\nu_3-\nu_2\unsure{-|+}3\nu_2\nu_3+3\nu_0]}_{\nu'}
\]
and \unsure{by|for} \unsure{congruence|congruences} modulo $6$
\[
\tag{87}
\mu=1-2\nu_3+6\underbrace{[\nu_3-\nu_2-2\nu_2\nu_3+2\nu_0]}_{\nu''},
\]
where the underline in $\underline{\mu}$ denotes the section in $\mathbb O_{12}^*$, the non-underlined \inserted{$\mu$} \deleted{\unsure{a}} \inserted{the} \inserted{\unsure{corresponding|corresponding}} section of $\mathbf G_m=\mathbb O_1^*$. We will indeed have
\[
\begin{aligned}
2\nu'-\nu_2&=2\nu_3-3\nu_2-6\nu_2\nu_3+6\nu_0\qquad\unsure{''}\\
3\nu''-\nu_3&=2\nu_3-3\nu_2-6\nu_2\nu_3+6\nu_0
\end{aligned}
\]

In summary

\uline{Proposition} Let us introduce on \inserted{the smooth \unsure{group|group} scheme} $\mathbb O_{12}^*$ the functions \inserted{$\underline{\mu}\mapsto$} $\unsure{\nu_2|\xi_2},\nu_3\unsure{,}$ \deleted{$\unsure{\xi_3|\varepsilon_3}$ \unsure{and}} (with values in $\{0,1\}_{S_0}\subset\inserted{\mathbf G_a=}\mathbb O_1$), $\nu',\nu''$ \deleted{$\unsure{\mu_3|\nu_3}$} (with values in $\unsure{\mathbb O_1|\mathbf G_a}$) and $\mu$ (with values in $\mathbf G_m\inserted{^{\unsure{*}}=\mathbb O_1^*\subset\mathbb O_1}$) \unsure{defined} by the conditions
\[
\tag{88}
\left\{\begin{array}{ll}
\mu={}&\text{canonical \unsure{image} of }\underline{\mu}\text{ in }\mathbb O_1^*\text{ under}\\
&\text{the homomorphism of Groups }\mathbb O_N^*\unsure{\longrightarrow|\mathrel{-}}\mathbb O_1^*\\[4pt]
\begin{gathered}
\inserted{(-1)^{\nu_2}}\\[-2pt]
\deleted{\text{\unsure{the image}}}\\[-2pt]
\deleted{(\unsure{\varepsilon_2|\xi_2})}
\end{gathered}={}&\text{canonical \unsure{image} of }\underline{\mu}\text{ in }(\mathbb Z/4\mathbb Z)^*_{S_0}\\[4pt]
(-1)^{\nu_3}={}&\rule{5cm}{0.4pt}\quad(\mathbb Z/6\mathbb Z)^*_{S_0}
\end{array}\right.
\]
whence
\[
\tag{89}
\deleted{\unsure{\Bigg\{}}\quad
\left\{\begin{aligned}
\underline{\mu}&=\deleted{\unsure{\mu_0}}(-5)^{\nu_2}(+5)^{\nu_3}+12[\unsure{\nu_0|\lambda_0}]\\
&=(1+4\nu_3)(1-6\nu_2)+12\nu_0\\
&=1+4\nu_3-6\nu_2+12\unsure{[}\nu_0-2\nu_2\nu_3]
\end{aligned}\right.
\]
for a well-determined section $\nu_0$ of $\mathbf G_a=\mathbb O_1$

\margin{56}
\margin{56\\\texttt{https://grothendieck.umontpellier.fr}\\Université de Montpellier, UMR CNRS 5149, Case courrier 051, Place Eugène Bataillon\\34095 Montpellier cedex 5 - France}